\section*{الفصل \ref{ch:b1:quantum-numbers} --- الأعداد الكمية والتوزيعات الإلكترونية}

\begin{solution}{exo:b1:quantum-numbers:1}
(أ) مستحيلة: $l \le n - 1 = 1$. (ب) ممكنة (إلكترون 3p).
(ج) مستحيلة: $m_s = \pm\tfrac12$. (د) ممكنة (إلكترون 4f).
(ه) مستحيلة: من أجل $l = 0$، لا يأخذ $m_l$ إلا القيمة 0.
\end{solution}

\begin{solution}{exo:b1:quantum-numbers:2}
3d: 5 مدارات، و10 إلكترونات؛ 4f: 7 مدارات، و14 إلكترونًا؛ 5p: 3
مدارات، و6 إلكترونات؛ الطبقة $n = 2$: $n^2 = 4$ مدارات، و$2n^2 = 8$
إلكترونات.
\end{solution}

\begin{solution}{exo:b1:quantum-numbers:3}
\ce{O}: $1\text{s}^2 2\text{s}^2 2\text{p}^4 = [\ce{He}]\,2\text{s}^2 2\text{p}^4$،
6 \omterm{def:b1:quantum-numbers:configuration}{إلكترونات تكافؤ}. \ce{P}: $1\text{s}^2 2\text{s}^2 2\text{p}^6 3\text{s}^2
3\text{p}^3 = [\ce{Ne}]\,3\text{s}^2 3\text{p}^3$، 5. \ce{K}: $[\ce{Ar}]\,4\text{s}^1$،
1. \ce{Ca}: $[\ce{Ar}]\,4\text{s}^2$، 2. \ce{Br}: $[\ce{Ar}]\,3\text{d}^{10}
4\text{s}^2 4\text{p}^5$، 7 (فالطبقة 3d الممتلئة جزء من القلب).
% ledger: gs:K, gs:Ca
\end{solution}

\begin{solution}{exo:b1:quantum-numbers:4}
\ce{N}، 2p: \omorbs{u,u,u}، 3 منفردة. \ce{S}، 3p: \omorbs{ud,u,u}، 2
منفردان. \ce{Cl}، 3p: \omorbs{ud,ud,u}، 1 منفرد. وفي كل حالة تكون 2s أو
3s ممتلئة، \omorbs{ud}.
\end{solution}

\begin{solution}{exo:b1:quantum-numbers:5}
\ce{Na+} و\ce{Mg^{2+}} و\ce{Al^{3+}} و\ce{F-} و\ce{O^{2-}}: $1\text{s}^2
2\text{s}^2 2\text{p}^6$، متساوية إلكترونيًا مع النيون. \ce{Cl-} و\ce{S^{2-}}
و\ce{Ca^{2+}}: $[\ce{Ne}]\,3\text{s}^2 3\text{p}^6$، متساوية إلكترونيًا مع
الأرغون.
\end{solution}

\begin{solution}{exo:b1:quantum-numbers:6}
تغادر إلكترونات 4s أولًا (\cref{prop:b1:quantum-numbers:ions}).
\ce{Ti^{2+}}: $[\ce{Ar}]\,3\text{d}^2$، 2 منفردان؛ \ce{Mn^{2+}}:
$3\text{d}^5$، 5؛ \ce{Co^{2+}}: $3\text{d}^7$، 3 (\omorbs{ud,ud,u,u,u})؛
\ce{Ni^{2+}}: $3\text{d}^8$، 2؛ \ce{Zn^{2+}}: $3\text{d}^{10}$، 0.
% ledger: gs:Ti2+, gs:Mn2+, gs:Co2+, gs:Zn2+
\end{solution}

\begin{solution}{exo:b1:quantum-numbers:7}
$\Delta E = 13.6\,(1/9 - 1/16) = 13.6 \times 7/144 = \qty{0.661}{eV}$؛
$\lambda = 1240/0.661 = \qty{1876}{nm} \approx \qty{1.88}{\micro m}$:
تحت الأحمر (أول خط في سلسلة باشن).
\end{solution}

\begin{solution}{exo:b1:quantum-numbers:8}
(أ) $[\ce{Ne}]\,3\text{s}^2 3\text{p}^4$: الكبريت، $Z = 16$.
(ب) $[\ce{Ar}]\,3\text{d}^{10} 4\text{s}^2 4\text{p}^5$: البروم، $Z = 35$.
(ج) $[\ce{Kr}]\,5\text{s}^1$: الروبيديوم، $Z = 37$.
(د) $[\ce{Ar}]\,3\text{d}^3 4\text{s}^2$: الفاناديوم، $Z = 23$.
إلكترون 3p للكبريت: $(3, 1, -1, +\tfrac12)$ مثلًا.
% ledger: gs:V
\end{solution}

\begin{solution}{exo:b1:quantum-numbers:9}
من أجل 4s، $n + l = 4 + 0 = 4$؛ ومن أجل 3d، $n + l = 3 + 2 = 5$: فتأتي 4s
أولًا. وبعد \babelsublr{3p ($n + l = 4$)} تكون \omterm{def:b1:quantum-numbers:orbital}{الطبقة الفرعية} التالية 4s ($n + l = 4$، و$n$
أكبر)، وهي التي تفتح الدورة 4؛ ولا تأتي \babelsublr{3d ($n + l = 5$)} إلا بعد 4s.
فلا تحوي الدورة 3 إلا 3s و3p: $2 + 6 = 8$ عناصر.
\end{solution}

\begin{solution}{exo:b1:quantum-numbers:10}
طاقة التأين $Z^2E_R = 4 \times 13.6 = \qty{54.4}{eV}$، أي أربعة أضعاف
طاقة تأين الهيدروجين. $2 \to 1$: $\Delta E = 54.4\,(1 - 1/4) = \qty{40.8}{eV}$،
$\lambda = 1240/40.8 = \qty{30.4}{nm}$، أي أقصر بأربع مرات من خط
الهيدروجين عند \qty{121.6}{nm}: فكل طاقة مضروبة في $Z^2 = 4$.
\end{solution}

\begin{solution}{exo:b1:quantum-numbers:11}
(أ) \omterm{def:b1:quantum-numbers:energy-level}{حالة مثارة} للبيريليوم (حالته الأساسية $1\text{s}^2 2\text{s}^2$).
(ب) مستحيل: لا تتسع 2p لأكثر من 6 إلكترونات. (ج) \omterm{def:b1:quantum-numbers:energy-level}{حالة مثارة}
للبيريليوم. (د) \omterm{def:b1:quantum-numbers:energy-level}{الحالة الأساسية} للفوسفور. (ه) مستحيل: \omterm{def:b1:quantum-numbers:orbital}{الطبقة الفرعية} s
تتسع لإلكترونين (باولي).
\end{solution}

\begin{solution}{exo:b1:quantum-numbers:12}
\ce{Fe^{3+}} $3\text{d}^5$: 5؛ \ce{Mn^{2+}} $3\text{d}^5$: 5؛ \ce{Fe^{2+}}
$3\text{d}^6$: 4؛ \ce{Cu^{2+}} $3\text{d}^9$: 1؛ \ce{Zn^{2+}}
$3\text{d}^{10}$: 0. الترتيب: $\ce{Fe^{3+}} = \ce{Mn^{2+}} > \ce{Fe^{2+}}
> \ce{Cu^{2+}} > \ce{Zn^{2+}}$. للحديد(III) \omterm{def:b1:quantum-numbers:unpaired}{إلكترون منفرد} زائد على
الحديد(II)، فينجذب أكثر؛ وليس للزنك(II) أي \omterm{def:b1:quantum-numbers:unpaired}{إلكترون منفرد}، فلا
ينجذب.
% ledger: gs:Fe2+, gs:Fe3+, gs:Mn2+, gs:Cu2+, gs:Zn2+
\end{solution}

\begin{solution}{pb:b1:quantum-numbers:1}
\textbf{1.} $E_1 = \qty{-13.6}{eV}$، $E_2 = \qty{-3.40}{eV}$،
$E_3 = \qty{-1.51}{eV}$، $E_4 = \qty{-0.85}{eV}$.
\textbf{2.} $\Delta E = 3.40 - 1.51 = \qty{1.89}{eV}$، $\lambda = 1240/1.89
= \qty{656}{nm}$: أحمر.
\textbf{3.} $4 \to 2$: \qty{2.55}{eV}، \qty{486}{nm} (أزرق مخضرّ)؛
$5 \to 2$: \qty{2.86}{eV}، \qty{434}{nm} (بنفسجي)؛ $6 \to 2$: \qty{3.022}{eV}،
\qty{410}{nm} (بنفسجي).
\textbf{4.} $7 \to 2$: \qty{3.12}{eV}، \qty{397}{nm}، أي في فوق البنفسجي؛
والانتقالات الأعلى أقصر موجةً أيضًا، و$3 \to 2$ هو أطول طول موجة
في سلسلة بالمر. وخطوط ليمان كلها في فوق البنفسجي وخطوط باشن
كلها في تحت الأحمر، فلا يُرى إلا أربعة خطوط.
\textbf{5.} $\infty \to 2$: \qty{3.40}{eV}، $\lambda = \qty{365}{nm}$.
\textbf{6.} $\Delta E = 13.6 \times 3/4 = \qty{10.2}{eV}$، $\lambda =
\qty{122}{nm}$: فوق بنفسجي بعيد، لا تراه العين ويمتصه
الزجاج والهواء.
\textbf{7.} يجب أن يحمل الفوتون \qty{13.6}{eV} على الأقل: $\lambda \le
1240/13.6 = \qty{91.2}{nm}$.
\textbf{8.} $l = 0$ ($m_l = 0$)، $l = 1$ ($m_l = -1, 0, 1$)، $l = 2$
($m_l = -2, \dots, 2$): $1 + 3 + 5 = 9$ مدارات.
\textbf{9.} $\sum_{l=0}^{n-1}(2l+1) = n^2$ مدارًا، في كل منها إلكترونان:
$2n^2$؛ ومن أجل $n = 4$، 32.
\textbf{10.} 1s، 2s، 2p، 3s، 3p، 4s، 3d، 4p، 5s، 4d، 5p، 6s، 4f، 5d، 6p،
7s، 5f، 6d، 7p.
\textbf{11.} تمتد كل دورة من \babelsublr{$n$s} إلى \babelsublr{$n$p}: 2، 8، 8، 18، 18، 32، 32.
\textbf{12.} تأتي \babelsublr{3d ($n + l = 5$)} بعد \babelsublr{4s ($n + l = 4$)}، التي تفتح
الدورة 4.
\textbf{13.} $Z = 2, 10, 18, 36, 54, 86, 118$.
\textbf{14.} $[\ce{Ar}]\,3\text{d}^6 4\text{s}^2$؛ 3d: \omorbs{ud,u,u,u,u}،
4 إلكترونات منفردة.
\textbf{15.} \ce{Fe^{2+}} $[\ce{Ar}]\,3\text{d}^6$: 4 منفردة؛
\ce{Fe^{3+}} $[\ce{Ar}]\,3\text{d}^5$: 5 منفردة، \omterm{def:b1:quantum-numbers:orbital}{وطبقة فرعية} نصف ممتلئة.
\textbf{16.} المتوقَّع $3\text{d}^4 4\text{s}^2$: 4 منفردة؛ والمرصود
$3\text{d}^5 4\text{s}^1$: 6 منفردة (5 في 3d، و1 في 4s).
\textbf{17.} \ce{Cu} $[\ce{Ar}]\,3\text{d}^{10} 4\text{s}^1$، \ce{Cu+}
$[\ce{Ar}]\,3\text{d}^{10}$، \ce{Cu^{2+}} $[\ce{Ar}]\,3\text{d}^9$.
% ledger: gs:Cu, gs:Cu+, gs:Cu2+
\textbf{18.} \ce{Cu^{2+}}، \omterm{def:b1:quantum-numbers:unpaired}{بإلكترون منفرد} واحد؛ وليس للأيون \ce{Cu+} أي \omterm{def:b1:quantum-numbers:unpaired}{إلكترون منفرد}.
\textbf{19.} \ce{Sc^{3+}} $[\ce{Ar}]$: 0؛ \ce{Ti^{3+}} $[\ce{Ar}]\,3\text{d}^1$:
1؛ \ce{Mn^{2+}} $[\ce{Ar}]\,3\text{d}^5$: 5. فللأيون \ce{Mn^{2+}} أكبر عدد.
\textbf{20.} ثلاثة إلكترونات في كل مدار؛ و$3(2l + 1)$ في كل \omterm{def:b1:quantum-numbers:orbital}{طبقة فرعية}؛
و$3n^2$ في كل طبقة.
\textbf{21.} 1s: 3؛ 2s: 3؛ 2p: 9؛ 3s: 3؛ 3p: 9.
\textbf{22.} تُغلق الدورة الأولى عندما تمتلئ 1s: $Z = 3$.
\textbf{23.} $Z = 4$، أي إلكترون واحد في 2s زائد على القلب النبيل.
\textbf{24.} 2s و2p: $3 + 9 = 12$ عنصرًا.
\textbf{25.} $3 + 12 = 15$: للغاز النبيل الثاني في عالم المغازل الثلاثة
العدد الذري \textbf{$Z = 15$}.
\end{solution}
